\chapter{Testes de Hipótese}

No capítulo anterior, estimamos parâmetros populacionais --- uma média ou uma proporção --- e expressamos a incerteza dessa estimativa através de um intervalo de confiança. O teste de hipóteses aborda um problema relacionado, porém distinto: em vez de perguntar "qual é o valor do parâmetro?", perguntamos "os dados são compatíveis com um valor específico (ou com a ausência de um efeito) que suponho a priori?"

\section{Como pensar antes de escolher um teste}

Uma dificuldade comum no início do estudo da estatística é imaginar que cada
pergunta científica corresponde a um teste estatístico específico. Essa
abordagem transforma a estatística em uma coleção de procedimentos a serem
memorizados: teste-$t$, qui-quadrado, Fisher, ANOVA, correlação e assim por
diante.

Uma abordagem mais útil é começar pela \textbf{pergunta científica} e, a
partir dela, determinar qual método estatístico é apropriado. A escolha do
teste não deve ser feita simplesmente pelo nome da variável ou por uma regra
isolada, mas pelo conjunto de características do problema.

Uma forma de organizar esse raciocínio é:

\begin{enumerate}
\item Pergunta científica
\item Desenho do estudo
\item Tipos de variáveis e estrutura dos grupos
\item Pressupostos
\item Método estatístico
\end{enumerate}

A \textbf{pergunta científica} determina o que desejamos investigar. Podemos
querer descrever uma população, comparar grupos, investigar uma associação,
estimar um efeito ou avaliar uma hipótese causal.

O \textbf{desenho do estudo} fornece informações fundamentais para essa
escolha. Dados provenientes de um ensaio clínico randomizado, por exemplo,
possuem uma estrutura diferente daquela encontrada em um estudo observacional.
Da mesma forma, medidas repetidas no mesmo indivíduo não podem ser tratadas
como se fossem observações independentes.

Em seguida, devemos considerar o \textbf{tipo de variável} envolvida. Uma variável quantitativa, como idade ou pressão arterial, exige procedimentos diferentes daqueles utilizados para uma variável categórica, como sexo, tratamento ou ocorrência de um evento.

Também é necessário observar a \textbf{estrutura dos grupos}. Quando existem
dois grupos, por exemplo, devemos perguntar se eles são independentes ou se as observações estão pareadas. Quando existem três ou mais grupos, a estratégia de análise pode ser diferente. A mesma variável pode, portanto, ser analisada por métodos distintos dependendo da maneira como os dados foram coletados.

Por fim, devem ser avaliados os \textbf{pressupostos} do procedimento. Alguns
métodos dependem de condições relacionadas à distribuição dos dados,
independência das observações, homogeneidade das variâncias ou tamanho das
contagens esperadas. Quando essas condições não são adequadas, outro
procedimento pode ser mais apropriado.

Assim, o teste estatístico deve ser visto como a \textbf{consequência de uma
sequência de decisões}, e não como o ponto de partida da análise.

Neste capítulo, serão apresentados os principais conceitos que fundamentam os
testes de hipótese. No capítulo seguinte, esses procedimentos serão organizados de forma prática, mostrando \textbf{quando utilizar cada um} e como escolher entre alternativas possíveis diante de diferentes tipos de dados e desenhos de estudo.

\section{Associação, diferença e causalidade}

Duas variáveis estão associadas quando a distribuição de uma varia de acordo
com os valores ou categorias da outra.Por exemplo, pode existir associação entre idade e pressão arterial se os
valores de pressão arterial tenderem a ser diferentes em indivíduos de
idades diferentes.

Uma comparação entre grupos também pode ser vista como uma forma particular
de estudar associação. Ao comparar a idade entre os grupos RF e Cryo, estamos
investigando se a distribuição da idade varia de acordo com o grupo de
tratamento.

É importante, entretanto, distinguir três perguntas:

\begin{itemize}
    \item \textbf{Descrição:} quais são os valores observados nos dados?
    \item \textbf{Associação:} existe relação estatística entre as variáveis?
    \item \textbf{Causalidade:} uma variável produz mudança na outra?
\end{itemize}

Essas perguntas não são equivalentes.

Associação não implica causalidade. Uma associação observada pode decorrer de:

\begin{itemize}
    \item efeito causal;
    \item causalidade reversa;
    \item confundimento;
    \item viés de seleção;
    \item viés de mensuração;
    \item acaso.
\end{itemize}

Em um ensaio clínico randomizado, a randomização ajuda a equilibrar, em
expectativa, fatores de prognóstico entre os grupos. Quando o estudo é
adequadamente conduzido, isso fortalece a interpretação causal da diferença
entre os tratamentos.

Em um estudo observacional, entretanto, uma associação entre exposição e
desfecho não estabelece, por si só, que a exposição causou o desfecho.

\section{Hipótese Nula e Alternativa}
Todo teste de hipóteses parte de duas afirmações mutuamente excludentes sobre a população:
\begin{itemize}
    \item \textbf{Hipótese nula ($H_0$):} afirma que não há efeito, diferença ou associação --- é o estado de "nada de novo", tomado como verdadeiro até que haja evidência suficiente em contrário. Ex.: \emph{"o novo medicamento tem o mesmo efeito que o placebo"}.
    \item \textbf{Hipótese alternativa ($H_1$ ou $H_a$):} afirma exatamente o oposto --- existe efeito, diferença ou associação. Ex.: \emph{"o novo medicamento reduz a pressão arterial mais do que o placebo"}.
\end{itemize}

A lógica é semelhante à de um julgamento: presume-se a inocência ($H_0$) e só se condena (rejeita-se $H_0$) diante de evidência que torne essa hipótese muito improvável. O teste nunca "prova" $H_0$ verdadeira --- apenas indica se há ou não evidência suficiente para rejeitá-la.

\begin{error}
\boxtitle{Direcionalidade e p-hacking}
A hipótese alternativa pode ser \textbf{bilateral} ($H_1: \mu \neq \mu_0$, o efeito pode ir em qualquer direção) ou \textbf{unilateral} ($H_1: \mu > \mu_0$ ou $H_1: \mu < \mu_0$, quando só uma direção do efeito é biologicamente plausível ou relevante). A escolha deve ser feita antes de olhar os dados --- decidir isso depois de ver o resultado é uma forma de viés conhecida como \emph{p-hacking}.
\end{error}

\section{O Valor-p: Definição correta e interpretação clínica}

\begin{definition}[Definição: Valor-p]
\textbf{Valor-p:} A probabilidade de observar um resultado tão extremo quanto (ou mais extremo do que) o obtido na amostra, assumindo que a hipótese nula é verdadeira.
\end{definition}

Um valor-$p$ pequeno indica que o resultado observado seria raro se $H_0$ fosse de fato verdadeira --- o que enfraquece a credibilidade de $H_0$. Convencionalmente, compara-se o valor-$p$ a um limiar de significância ($\alpha$), geralmente $0{,}05$:
\begin{itemize}
    \item $p < \alpha$: rejeita-se $H_0$ --- resultado "estatisticamente significativo".
    \item $p \ge \alpha$: não se rejeita $H_0$ --- não há evidência suficiente para afastar a hipótese de ausência de efeito.
\end{itemize}

\begin{deepdive}
\boxtitle{Aprofundando: O que o valor-p NÃO significa}

\begin{itemize}
    \item Não é a probabilidade de $H_0$ ser verdadeira, nem de $H_1$ ser verdadeira.
    \item Não mede o tamanho ou a relevância clínica do efeito --- uma amostra muito grande pode gerar $p < 0{,}05$ para uma diferença clinicamente irrelevante (ex.: $0{,}1\text{ mmHg}$ de diferença de pressão arterial).
    \item Não é a probabilidade de o achado ser "por acaso" no sentido coloquial, e $p = 0{,}06$ não é qualitativamente diferente de $p = 0{,}04$ --- o limiar de $0{,}05$ é uma convenção, não uma lei natural.
\end{itemize}
Por isso, diretrizes editoriais recomendam reportar sempre o intervalo de confiança e uma medida de tamanho de efeito ao lado do valor-$p$.
\end{deepdive}

\section{Exemplo em R}

Suponha que, em uma população de referência, a pressão arterial sistólica
média seja de $130$ mmHg. Um pesquisador deseja investigar se pacientes
submetidos a determinado tratamento apresentam uma pressão arterial sistólica
média diferente desse valor.

Uma amostra de $100$ pacientes apresentou média de $134$ mmHg. Para simplificar
o exemplo, suponha que o desvio-padrão populacional seja conhecido e igual a
$20$ mmHg.

A pergunta científica é, portanto:

\begin{quote}
A pressão arterial sistólica média dos pacientes é diferente de
$130$ mmHg?
\end{quote}

Como a pergunta é sobre uma única média populacional e o desvio-padrão
populacional é conhecido, podemos utilizar um teste $Z$ para uma média.

\subsection{Hipóteses}

O primeiro passo é transformar a pergunta científica em hipóteses
estatísticas.

A hipótese nula representa o valor de referência:

\[
H_0:\mu=130.
\]

A hipótese alternativa representa a possibilidade de uma diferença:

\[
H_1:\mu\neq130.
\]

O teste é, portanto, \textbf{bilateral}, pois estamos interessados em detectar
uma diferença em qualquer direção.

\begin{itemize}
    \item $H_0$: a média populacional é $130$ mmHg;
    \item $H_1$: a média populacional é diferente de $130$ mmHg.
\end{itemize}

O nível de significância será definido como:

\[
\alpha=0{,}05.
\]

Isso significa que, sob as condições do teste, aceitamos uma probabilidade de
$5\%$ de cometer um erro Tipo I caso rejeitemos uma hipótese nula verdadeira.

\subsection{Calculando a estatística de teste}

Para o teste $Z$ de uma média, utilizamos:

\[
Z=
\frac{\bar{x}-\mu_0}
{\sigma/\sqrt{n}},
\]

em que $\bar{x}$ é a média observada na amostra, $\mu_0$ é o valor especificado
pela hipótese nula, $\sigma$ é o desvio-padrão populacional e $n$ é o tamanho
da amostra.

Neste exemplo:

\[
\bar{x}=134,\qquad
\mu_0=130,\qquad
\sigma=20,\qquad
n=100.
\]

Portanto:

\[
Z=
\frac{134-130}{20/\sqrt{100}}
=
\frac{4}{2}
=
2.
\]

A estatística de teste é, portanto, $Z=2$.

Podemos realizar esse cálculo diretamente em R.

\begin{code}[Código: cálculo da estatística Z]
# Média observada na amostra
media <- 134

# Valor da média especificado em H0
media_nula <- 130

# Desvio-padrão populacional conhecido
dp_populacional <- 20

# Tamanho da amostra
n <- 100

# Cálculo da estatística Z
z <- (media - media_nula) /
  (dp_populacional / sqrt(n))

z
\end{code}

Cada parte do código corresponde diretamente aos elementos da fórmula
estatística:

\begin{itemize}
    \item \texttt{media} representa a média observada, $\bar{x}$;
    \item \texttt{media\_nula} representa o valor proposto por $H_0$, $\mu_0$;
    \item \texttt{dp\_populacional} representa $\sigma$;
    \item \texttt{n} representa o tamanho da amostra;
    \item a última expressão calcula a estatística $Z$.
\end{itemize}

O resultado é:

\begin{code}[Resultado da estatística Z]
[1] 2
\end{code}

Isso significa que a média observada está a $2$ erros-padrão do valor
especificado pela hipótese nula.

\subsection{Calculando o valor-p}

Como nossa hipótese alternativa é bilateral, devemos considerar resultados extremos tanto acima quanto abaixo de
$130$ mmHg. A distribuição utilizada é a distribuição normal padrão. Em R, a função
\texttt{pnorm()} fornece probabilidades acumuladas dessa distribuição.

Podemos calcular o valor-$p$ da seguinte maneira:

\begin{code}[Código: cálculo do valor-p]
# Probabilidade de observar um valor
# pelo menos tão extremo quanto |Z| = 2
p_valor <- 2 * (1 - pnorm(abs(z)))

p_valor
\end{code}

A função \texttt{abs()} transforma $Z=2$ em seu valor absoluto. Isso é útil
porque, em um teste bilateral, estamos interessados na distância da estatística
em relação a zero, independentemente da direção.

A expressão $1-\texttt{pnorm}(|Z|)$ calcula a probabilidade de observar um valor pelo menos tão grande quanto o
valor observado em uma das caudas. Multiplicamos essa probabilidade por $2$
porque o teste é bilateral.

O resultado é aproximadamente:

\begin{code}[Resultado do valor-p]
[1] 0.0455
\end{code}

Portanto,

\[
p\approx0{,}0455.
\]

\subsection{Tomando a decisão}

Agora podemos comparar o valor-$p$ com o nível de significância definido
anteriormente:

\[
p=0{,}0455 < \alpha=0{,}05.
\]

Assim, rejeitamos $H_0$.

Podemos também fazer essa decisão diretamente no R:

\begin{code}[Código: decisão do teste]
# Nível de significância
alpha <- 0.05

# Regra de decisão
if (p_valor < alpha) {
  print("Rejeitar H0")
} else {
  print("Não rejeitar H0")
}
\end{code}

O resultado será:

\begin{code}[Resultado da decisão]
[1] "Rejeitar H0"
\end{code}

\subsection{Interpretação}

Podemos concluir que, neste exemplo, há \textbf{evidência estatística de que a
média populacional da pressão arterial sistólica seja diferente de
$130$ mmHg}.

O teste não demonstra que $H_1$ é verdadeira. Ele indica que os dados observados
seriam relativamente pouco compatíveis com $H_0$, assumindo que as condições
do teste sejam válidas. Também não podemos interpretar $p=0{,}0455$ como ``uma probabilidade de
4,55\% de $H_0$ ser verdadeira''. O valor-$p$ é calculado \textbf{assumindo
que $H_0$ é verdadeira}.

Além disso, a significância estatística não informa, por si só, se uma
diferença de $4$ mmHg é clinicamente relevante. Para essa avaliação, devemos
considerar a magnitude do efeito, sua incerteza e o contexto clínico.

\begin{error}
\boxtitle{Não rejeitar $H_0$ não significa provar que não existe efeito}

Se o valor-$p$ fosse maior que $0{,}05$, a conclusão seria ``não rejeitar
$H_0$''. Isso não significaria que a média populacional é necessariamente
igual a $130$ mmHg.

Na verdade, significaria apenas que os dados não forneceram evidência estatística
suficiente para rejeitar esse valor de referência, considerando o procedimento
e o nível de significância adotados.

Essa distinção é importante porque a capacidade de detectar uma diferença
depende, entre outros fatores, do tamanho da amostra, da variabilidade dos
dados e do tamanho do efeito. Esses fatores serão discutidos posteriormente
ao estudarmos o poder estatístico.
\end{error}

\section{Erros de Decisão}
Como toda decisão baseada em amostra, o teste de hipóteses está sujeito a erro. Existem dois tipos possíveis, resumidos na Tabela~\ref{tab:erros_decisao}.

\begin{center}
\captionof{table}{Matriz de Decisão e Tipos de Erro em Testes de Hipóteses}
\label{tab:erros_decisao}

\vspace{2mm}

\begin{tabular}{l c c}
\toprule
\textbf{Decisão do Teste} & \textbf{$H_0$ Verdadeira} & \textbf{$H_0$ Falsa} \\
\midrule
\textbf{Rejeita $H_0$} &
Erro Tipo I ($\alpha$) --- Falso Positivo &
Decisão Correta (Poder = $1-\beta$) \\

\textbf{Não rejeita $H_0$} &
Decisão Correta &
Erro Tipo II ($\beta$) --- Falso Negativo \\
\bottomrule
\end{tabular}
\end{center}

\subsection{Erro Tipo I}
Ocorre quando rejeitamos $H_0$ sendo ela, na realidade, verdadeira --- concluímos que há um efeito que de fato não existe (falso positivo). A probabilidade máxima aceita para esse erro é o próprio nível de significância $\alpha$ (tipicamente $0{,}05$, ou seja, $5\%$ de chance).

\subsection{Erro Tipo II}
Ocorre quando não rejeitamos $H_0$ sendo ela, na realidade, falsa --- deixamos de detectar um efeito que de fato existe (falso negativo). Sua probabilidade é denotada por $\beta$, e o complemento ($1-\beta$) é chamado de \textbf{poder do teste}.
\begin{error}
\boxtitle{Confundir Significância Estatística com Relevância Clínica}

Amostras excessivamente grandes aumentam o poder estatístico a ponto de produzir $p < 0{,}05$ para diferenças minúsculas (por exemplo, uma redução de $0{,}1\text{ mmHg}$ na pressão arterial) sem qualquer impacto prático no prognóstico do paciente.
\end{error}

\section{Poder estatístico}

O \textbf{poder estatístico} está diretamente relacionado ao Erro Tipo II. Se
$\beta$ é a probabilidade de não rejeitar $H_0$ quando existe um efeito de
determinado tamanho, então:

\[
\text{Poder}=1-\beta.
\]

Em outras palavras, o poder é a probabilidade de o teste \textbf{rejeitar
corretamente $H_0$ quando uma diferença ou efeito de determinado tamanho
realmente existe}.

Isso significa que o poder não é uma propriedade fixa de um teste ou de um
estudo. Ele depende do tamanho do efeito que se deseja detectar. Um estudo
pode ter alto poder para detectar uma diferença grande, mas poder muito menor
para detectar uma diferença pequena.

\begin{definition}[Poder estatístico]
\textbf{Poder estatístico:} probabilidade de rejeitar corretamente a hipótese
nula quando existe, na população, um efeito de determinado tamanho sob a
hipótese alternativa.

\[
\text{Poder}=1-\beta.
\]
\end{definition}

O poder é determinado principalmente por quatro componentes:

\begin{itemize}
    \item \textbf{tamanho da amostra ($n$):} amostras maiores geralmente
    aumentam o poder;

    \item \textbf{tamanho do efeito:} efeitos maiores são mais fáceis de
    detectar e, portanto, aumentam o poder;

    \item \textbf{nível de significância ($\alpha$):} aumentar $\alpha$
    torna mais fácil rejeitar $H_0$ e, consequentemente, aumenta o poder,
    embora também aumente a probabilidade de Erro Tipo I;

    \item \textbf{variabilidade dos dados:} quanto maior a variabilidade
    em relação ao efeito que se deseja detectar, mais difícil é identificar
    esse efeito e menor tende a ser o poder.
\end{itemize}

Podemos representar essa relação de forma simplificada como:

\[
\text{Poder}
=
f(n,\text{tamanho do efeito},\alpha,\text{variabilidade}).
\]

\subsection{Poder e tamanho da amostra}

Essa relação é particularmente importante no planejamento de estudos.

Suponha que um pesquisador queira comparar dois tratamentos e considere que
uma diferença de $5$ mmHg na pressão arterial seja clinicamente relevante.
Antes de coletar os dados, ele pode perguntar:

\begin{quote}
Quantos participantes são necessários para que o estudo tenha, por exemplo,
$90\%$ de probabilidade de detectar uma diferença de $5$ mmHg, caso ela
realmente exista?
\end{quote}

Essa é uma pergunta de \textbf{cálculo de tamanho amostral baseado em poder}.

Em um planejamento prospectivo, o pesquisador geralmente especifica:

\begin{itemize}
    \item o efeito mínimo considerado clinicamente relevante;
    \item a variabilidade esperada dos dados;
    \item o nível de significância $\alpha$;
    \item o poder desejado, frequentemente $80\%$ ou $90\%$.
\end{itemize}

A partir dessas informações, determina-se o tamanho de amostra necessário.

Por exemplo, se um estudo for planejado com poder de $90\%$, isso significa que,
\textbf{se o efeito especificado no planejamento realmente existir}, o
procedimento estatístico terá aproximadamente $90\%$ de probabilidade de
rejeitar $H_0$ nas condições assumidas.

Isso não significa que o estudo terá $90\%$ de probabilidade de encontrar
qualquer efeito que exista. Um efeito menor que aquele considerado no
planejamento pode ser muito mais difícil de detectar.

\begin{deepdive}
\boxtitle{Poder, tamanho do efeito e relevância clínica}

O conceito de poder ajuda a estabelecer uma conexão importante entre
estatística e medicina.

Considere dois estudos que utilizam o mesmo nível de significância e o mesmo
teste estatístico. O primeiro foi planejado para detectar uma diferença de
$10$ mmHg, enquanto o segundo foi planejado para detectar uma diferença de
$2$ mmHg.

Mesmo que tenham o mesmo número de participantes, o segundo estudo precisará
de uma amostra maior para alcançar o mesmo poder, pois detectar uma diferença
menor em relação à variabilidade dos dados é mais difícil.

Por isso, o tamanho amostral não deve ser escolhido apenas pela conveniência
ou pela disponibilidade de participantes. Idealmente, ele deve estar
relacionado a uma diferença que tenha significado científico ou clínico.

O conceito de \textbf{menor efeito clinicamente importante} é particularmente
útil nesse contexto: o objetivo não é simplesmente ter poder suficiente para
detectar algum efeito, mas ter poder suficiente para detectar um efeito que
seja relevante para a pergunta científica.
\end{deepdive}

\begin{error}
\boxtitle{Não confundir poder com probabilidade de o resultado ser significativo}
Um estudo com poder de $80\%$ não significa que exista uma probabilidade de
$80\%$ de obter $p<0{,}05$ independentemente do que acontecer nos dados.

O poder de $80\%$ é definido \textbf{para um determinado tamanho de efeito,
sob determinadas condições e pressupostos}. Se esse efeito realmente existir
e as demais condições do planejamento forem adequadas, a probabilidade de
rejeitar $H_0$ será aproximadamente $80\%$.

Da mesma forma, um estudo com baixo poder que produza $p\geq0{,}05$ não
permite concluir que o efeito seja inexistente. O resultado pode simplesmente
refletir uma capacidade insuficiente de detectar o efeito de interesse.
\end{error}

\begin{error}
\boxtitle{Poder pós-hoc}

O cálculo de poder é mais útil quando realizado \textbf{antes da coleta dos
dados}, durante o planejamento do estudo.

Depois que os dados já foram coletados e o resultado do teste é conhecido,
calcular um chamado ``poder observado'' com base no efeito estimado na própria
amostra geralmente acrescenta pouca informação. Esse cálculo está fortemente
relacionado ao valor-$p$ observado e pode ser enganoso como forma de explicar
um resultado não significativo.

Quando um estudo não encontra significância estatística, é mais informativo
examinar:

\begin{itemize}
    \item a estimativa do efeito;
    \item seu intervalo de confiança;
    \item a precisão da estimativa;
    \item o efeito mínimo clinicamente relevante;
    \item e, quando apropriado, o poder planejado do estudo.
\end{itemize}

Assim, a pergunta adequada não é simplesmente ``o estudo teve poder
suficiente?'', mas:

\begin{quote}
Os dados permitem excluir efeitos clinicamente relevantes com uma precisão
adequada?
\end{quote}
\end{error}

\subsection{Exemplo em R}

Quando um teste não encontra diferença significativa, surge uma pergunta
importante: será que realmente não há diferença, ou o estudo não tinha
capacidade suficiente para detectar um efeito clinicamente relevante?

\begin{code}[Cálculo prospectivo do tamanho amostral]
library(pwr)

# Diferença clinicamente relevante
diferenca <- 5

# Desvio-padrão esperado
dp <- 18

# Tamanho de efeito padronizado
d <- diferenca / dp

# Parâmetros do planejamento
alpha <- 0.05
poder_desejado <- 0.90

# Calcular tamanho de amostra por grupo
planejamento <- pwr.t.test(
  d = d,
  sig.level = alpha,
  power = poder_desejado,
  alternative = "two.sided",
  type = "two.sample"
)

planejamento
\end{code}

\begin{code}[Resultado do cálculo]
        Two-sample t test power calculation 

            n = 273.3163
            d = 0.2777778
    sig.level = 0.05
        power = 0.9
alternative = two.sided

NOTE: n is number in *each* group
\end{code}

O resultado informa o número aproximado de participantes necessários em cada
grupo para alcançar o poder especificado, dadas as demais premissas do
planejamento.

É importante perceber que o poder é definido para um \textbf{determinado
tamanho de efeito}. Um estudo pode ter alto poder para detectar uma diferença
grande e poder muito menor para detectar uma diferença pequena.

\begin{deepdive}
\boxtitle{Por que usar o cálculo baseado em poder?}

O cálculo baseado em poder é mais adequado quando o objetivo é planejar um
estudo para detectar uma \textbf{diferença clinicamente relevante}. Nesse caso, especificamos previamente o tamanho do efeito que desejamos detectar, a
variabilidade esperada, o nível de significância e o poder desejado.

A fórmula baseada na margem de erro, por outro lado, é apropriada quando o
objetivo principal é obter uma estimativa com determinada \textbf{precisão},
por exemplo, estimar uma proporção com margem de erro de $5\%$ em um intervalo
de confiança de $95\%$. Ela não considera diretamente qual diferença entre
grupos seria clinicamente importante nem a probabilidade de detectar essa
diferença.

Assim, para um estudo comparativo ou experimental, em que queremos saber
\emph{quantos participantes são necessários para detectar um efeito relevante}, o cálculo baseado em poder é geralmente a abordagem mais adequada.
\end{deepdive}


\section{Resumo}
\begin{itemize}
    \item \textbf{$H_0$ e $H_1$} representam hipóteses mutuamente excludentes
    sobre a população. A hipótese alternativa pode ser bilateral ou
    unilateral, e sua direção deve ser definida antes da análise dos dados.

    \item \textbf{O valor-$p$} quantifica a compatibilidade dos dados, ou de
    resultados ainda mais extremos, com a hipótese nula. Ele não representa a
    probabilidade de $H_0$ ser verdadeira.

    \item \textbf{Não rejeitar $H_0$ não significa provar sua verdade}.
    Ausência de evidência estatística suficiente para detectar uma diferença
    não demonstra que os grupos sejam iguais.

    \item \textbf{Erro Tipo I} ocorre quando rejeitamos uma hipótese nula
    verdadeira, enquanto o \textbf{Erro Tipo II} ocorre quando não rejeitamos
    uma hipótese nula falsa.

    \item \textbf{O poder estatístico} é a probabilidade de detectar um efeito
    de determinado tamanho quando esse efeito realmente existe. Ele depende,
    entre outros fatores, do tamanho da amostra, do tamanho do efeito e do
    nível de significância.

    \item \textbf{Significância estatística não é sinônimo de relevância
    clínica}. A interpretação de um resultado deve considerar a magnitude do
    efeito, seu intervalo de confiança e o contexto científico e clínico.

    \item \textbf{Associação não implica causalidade}. A interpretação causal
    depende do desenho do estudo e da possibilidade de controlar fontes de
    viés e confundimento.

    \item \textbf{A escolha do método estatístico começa pela pergunta
    científica}, passando pelo desenho do estudo, tipo de variável, estrutura
    dos grupos e pressupostos do procedimento.
\end{itemize}

